\par\noindent\SourceLeaf{52}%
\hspace*{\parindent}B) Quaeramus iam generatim legem coefficientis, quam cosinus multipli alicuius imparis v. g. $\cos(2q-1)\varphi$ nanciscitur. Inter terminos seriei $(1+n\cos\varphi)^{\nu}$ secundum potestates ipsius $\cos\varphi$ procedentis primus ille, qui $\cos\varphi^{2q-1}$ continet, secundum cosinus multiplorum evolutus, in termino suo $(q-1)^{\text{to}}$ multiplum requisitum continet. Est nimirum ex lege antea $(2,\beta)$ exposita, ipsius $\cos\varphi^{2q-1}$ terminus $(q-1)^{\text{tus}}$

\[
= \frac{1}{2^{2q-2}} {}^{2q-1}\mathfrak{B}^{(2q-1)} \cos(2q-1)\varphi,
\]

\noindent sive, cum sit ${}^{2q-1}\mathfrak{B}^{(2q-1)} = 1$,

\[
\frac{1}{2^{2q-2}} \cos(2q-1)\varphi..
\]

\noindent Erat autem in serie fundamentali terminus $\cos\varphi^{2q-1}$ continens ductus in ${}^{\nu}\mathfrak{B}^{(2q-1)} n^{2q-1}$. Quo factore adiecto habebitur pars prima coefficientis, qui ad $\cos(2q-1)\varphi$ pertinet, sive terminus initialis seriei, qua coefficiens ille exhibetur

\[
= \frac{1}{2^{2q-2}} {}^{\nu}\mathfrak{B}^{(2q-1)} n^{2q-1}.
\]

\noindent Iam vero quaelibet ipsius $\cos\varphi$ potestas impar, cuius exponens $2q-1$ excedit, si per cosinus multiplorum exprimatur, praebebit terminum in $\cos(2q-1)$ ductum; estque horum $h^{\text{tus}}$, qui ex $h^{\text{ta}}$ post $\cos\varphi^{2q-1}$ potestate impari, cuius igitur exponens est $2q+2h-1$, originem trahit. Sumitur autem ex quantitate

\[
{}^{\nu}\mathfrak{B}^{(2q+2h-1)} n^{2q+2h-1} \cos\varphi^{2q+2h-1}
\]

\noindent secundum cosinus multiplorum evoluta terminus $(q-1)^{\text{tus}}$; fit autem ille (cum sit ex lege cognita, ipsius $\cos\varphi^{2q+2h-1}$ terminus $(q-1)^{\text{tus}}$

\[
= \frac{1}{2^{2q+2h-2}} {}^{2q+2h-1}\mathfrak{B}^{(2q+h-1)} \cos(2q-1)) 
\]

\[
= \frac{1}{2^{2q+2h-2}} {}^{2q+2h-1}\mathfrak{B}^{(2q+h-1)}  .  {}^{\nu}\mathfrak{B}^{(2q+2h-1)} n^{2q+2h-1} \cos(2q-1)\varphi.
\]

\noindent Huius igitur quantitatis coefficiens est $h^{\text{tus}}$ illorum, quos ex successiva partium seriei fundamentalis evolutione adsumta $\cos(2q-1)\varphi$ sortitur. Quare si, uti fas est, in nova serie condenda istarum partium summam pro uno habeamus coefficiente, fiet ipsius $\cos(2q-1)\varphi$ coefficiens seriei aequalis, cuius terminus initialis

\[
= \frac{1}{2^{2q-2}} {}^{\nu}\mathfrak{B}^{(2q-1)} n^{2q-1}
\]
